% ========================================================
% BÀI TẬP TOÁN 5 - KẾT NỐI TRI THỨC VỚI CUỘC SỐNG
% Bài 54: Thực hành tính toán ước lượng thể tích hình khối
% Định dạng: BTTeX Part III - Trắc nghiệm điền số ngắn
% ========================================================

% Câu 1
\begin{ex}
  Một bể cá hình hộp chữ nhật có chiều dài $3\text{ dm}$, chiều rộng $2\text{ dm}$ và mực nước cao $1\text{ dm}$. Thể tích nước trong bể là bao nhiêu lít (biết $1\text{ dm}^3 = 1\text{ l}$)?
  \shortans{6}
  \loigiai{
    Thể tích là: 3 x 2 x 1 = 6 dm3 = 6 lít.
  }
\end{ex}

% Câu 2
\begin{ex}
  Một bể cá hình hộp chữ nhật có chiều dài $4\text{ dm}$, chiều rộng $2\text{ dm}$ và mực nước cao $1\text{ dm}$. Thể tích nước trong bể là bao nhiêu lít (biết $1\text{ dm}^3 = 1\text{ l}$)?
  \shortans{8}
  \loigiai{
    Thể tích là: 4 x 2 x 1 = 8 dm3 = 8 lít.
  }
\end{ex}

% Câu 3
\begin{ex}
  Một bể cá hình hộp chữ nhật có chiều dài $5\text{ dm}$, chiều rộng $2\text{ dm}$ và mực nước cao $1\text{ dm}$. Thể tích nước trong bể là bao nhiêu lít (biết $1\text{ dm}^3 = 1\text{ l}$)?
  \shortans{10}
  \loigiai{
    Thể tích là: 5 x 2 x 1 = 10 dm3 = 10 lít.
  }
\end{ex}

% Câu 4
\begin{ex}
  Một bể cá hình hộp chữ nhật có chiều dài $6\text{ dm}$, chiều rộng $2\text{ dm}$ và mực nước cao $1\text{ dm}$. Thể tích nước trong bể là bao nhiêu lít (biết $1\text{ dm}^3 = 1\text{ l}$)?
  \shortans{12}
  \loigiai{
    Thể tích là: 6 x 2 x 1 = 12 dm3 = 12 lít.
  }
\end{ex}

% Câu 5
\begin{ex}
  Một bể cá hình hộp chữ nhật có chiều dài $7\text{ dm}$, chiều rộng $2\text{ dm}$ và mực nước cao $1\text{ dm}$. Thể tích nước trong bể là bao nhiêu lít (biết $1\text{ dm}^3 = 1\text{ l}$)?
  \shortans{14}
  \loigiai{
    Thể tích là: 7 x 2 x 1 = 14 dm3 = 14 lít.
  }
\end{ex}

% Câu 6
\begin{ex}
  Một bể cá hình hộp chữ nhật có chiều dài $8\text{ dm}$, chiều rộng $2\text{ dm}$ và mực nước cao $1\text{ dm}$. Thể tích nước trong bể là bao nhiêu lít (biết $1\text{ dm}^3 = 1\text{ l}$)?
  \shortans{16}
  \loigiai{
    Thể tích là: 8 x 2 x 1 = 16 dm3 = 16 lít.
  }
\end{ex}

% Câu 7
\begin{ex}
  Một bể cá hình hộp chữ nhật có chiều dài $9\text{ dm}$, chiều rộng $2\text{ dm}$ và mực nước cao $1\text{ dm}$. Thể tích nước trong bể là bao nhiêu lít (biết $1\text{ dm}^3 = 1\text{ l}$)?
  \shortans{18}
  \loigiai{
    Thể tích là: 9 x 2 x 1 = 18 dm3 = 18 lít.
  }
\end{ex}

% Câu 8
\begin{ex}
  Một bể cá hình hộp chữ nhật có chiều dài $10\text{ dm}$, chiều rộng $2\text{ dm}$ và mực nước cao $1\text{ dm}$. Thể tích nước trong bể là bao nhiêu lít (biết $1\text{ dm}^3 = 1\text{ l}$)?
  \shortans{20}
  \loigiai{
    Thể tích là: 10 x 2 x 1 = 20 dm3 = 20 lít.
  }
\end{ex}

% Câu 9
\begin{ex}
  Một bể cá hình hộp chữ nhật có chiều dài $11\text{ dm}$, chiều rộng $2\text{ dm}$ và mực nước cao $1\text{ dm}$. Thể tích nước trong bể là bao nhiêu lít (biết $1\text{ dm}^3 = 1\text{ l}$)?
  \shortans{22}
  \loigiai{
    Thể tích là: 11 x 2 x 1 = 22 dm3 = 22 lít.
  }
\end{ex}

% Câu 10
\begin{ex}
  Một bể cá hình hộp chữ nhật có chiều dài $12\text{ dm}$, chiều rộng $2\text{ dm}$ và mực nước cao $1\text{ dm}$. Thể tích nước trong bể là bao nhiêu lít (biết $1\text{ dm}^3 = 1\text{ l}$)?
  \shortans{24}
  \loigiai{
    Thể tích là: 12 x 2 x 1 = 24 dm3 = 24 lít.
  }
\end{ex}

% Câu 11
\begin{ex}
  Một bể cá hình hộp chữ nhật có chiều dài $13\text{ dm}$, chiều rộng $2\text{ dm}$ và mực nước cao $1\text{ dm}$. Thể tích nước trong bể là bao nhiêu lít (biết $1\text{ dm}^3 = 1\text{ l}$)?
  \shortans{26}
  \loigiai{
    Thể tích là: 13 x 2 x 1 = 26 dm3 = 26 lít.
  }
\end{ex}

% Câu 12
\begin{ex}
  Một bể cá hình hộp chữ nhật có chiều dài $14\text{ dm}$, chiều rộng $2\text{ dm}$ và mực nước cao $1\text{ dm}$. Thể tích nước trong bể là bao nhiêu lít (biết $1\text{ dm}^3 = 1\text{ l}$)?
  \shortans{28}
  \loigiai{
    Thể tích là: 14 x 2 x 1 = 28 dm3 = 28 lít.
  }
\end{ex}

% Câu 13
\begin{ex}
  Một bể cá hình hộp chữ nhật có chiều dài $15\text{ dm}$, chiều rộng $2\text{ dm}$ và mực nước cao $1\text{ dm}$. Thể tích nước trong bể là bao nhiêu lít (biết $1\text{ dm}^3 = 1\text{ l}$)?
  \shortans{30}
  \loigiai{
    Thể tích là: 15 x 2 x 1 = 30 dm3 = 30 lít.
  }
\end{ex}

% Câu 14
\begin{ex}
  Một bể cá hình hộp chữ nhật có chiều dài $16\text{ dm}$, chiều rộng $2\text{ dm}$ và mực nước cao $1\text{ dm}$. Thể tích nước trong bể là bao nhiêu lít (biết $1\text{ dm}^3 = 1\text{ l}$)?
  \shortans{32}
  \loigiai{
    Thể tích là: 16 x 2 x 1 = 32 dm3 = 32 lít.
  }
\end{ex}

% Câu 15
\begin{ex}
  Một bể cá hình hộp chữ nhật có chiều dài $17\text{ dm}$, chiều rộng $2\text{ dm}$ và mực nước cao $1\text{ dm}$. Thể tích nước trong bể là bao nhiêu lít (biết $1\text{ dm}^3 = 1\text{ l}$)?
  \shortans{34}
  \loigiai{
    Thể tích là: 17 x 2 x 1 = 34 dm3 = 34 lít.
  }
\end{ex}

% Câu 16
\begin{ex}
  Một bể cá hình hộp chữ nhật có chiều dài $18\text{ dm}$, chiều rộng $2\text{ dm}$ và mực nước cao $1\text{ dm}$. Thể tích nước trong bể là bao nhiêu lít (biết $1\text{ dm}^3 = 1\text{ l}$)?
  \shortans{36}
  \loigiai{
    Thể tích là: 18 x 2 x 1 = 36 dm3 = 36 lít.
  }
\end{ex}

% Câu 17
\begin{ex}
  Một bể cá hình hộp chữ nhật có chiều dài $19\text{ dm}$, chiều rộng $2\text{ dm}$ và mực nước cao $1\text{ dm}$. Thể tích nước trong bể là bao nhiêu lít (biết $1\text{ dm}^3 = 1\text{ l}$)?
  \shortans{38}
  \loigiai{
    Thể tích là: 19 x 2 x 1 = 38 dm3 = 38 lít.
  }
\end{ex}

% Câu 18
\begin{ex}
  Một bể cá hình hộp chữ nhật có chiều dài $20\text{ dm}$, chiều rộng $2\text{ dm}$ và mực nước cao $1\text{ dm}$. Thể tích nước trong bể là bao nhiêu lít (biết $1\text{ dm}^3 = 1\text{ l}$)?
  \shortans{40}
  \loigiai{
    Thể tích là: 20 x 2 x 1 = 40 dm3 = 40 lít.
  }
\end{ex}

% Câu 19
\begin{ex}
  Một bể cá hình hộp chữ nhật có chiều dài $21\text{ dm}$, chiều rộng $2\text{ dm}$ và mực nước cao $1\text{ dm}$. Thể tích nước trong bể là bao nhiêu lít (biết $1\text{ dm}^3 = 1\text{ l}$)?
  \shortans{42}
  \loigiai{
    Thể tích là: 21 x 2 x 1 = 42 dm3 = 42 lít.
  }
\end{ex}

% Câu 20
\begin{ex}
  Một bể cá hình hộp chữ nhật có chiều dài $22\text{ dm}$, chiều rộng $2\text{ dm}$ và mực nước cao $1\text{ dm}$. Thể tích nước trong bể là bao nhiêu lít (biết $1\text{ dm}^3 = 1\text{ l}$)?
  \shortans{44}
  \loigiai{
    Thể tích là: 22 x 2 x 1 = 44 dm3 = 44 lít.
  }
\end{ex}

% Câu 21
\begin{ex}
  Một bể cá hình hộp chữ nhật có chiều dài $23\text{ dm}$, chiều rộng $2\text{ dm}$ và mực nước cao $1\text{ dm}$. Thể tích nước trong bể là bao nhiêu lít (biết $1\text{ dm}^3 = 1\text{ l}$)?
  \shortans{46}
  \loigiai{
    Thể tích là: 23 x 2 x 1 = 46 dm3 = 46 lít.
  }
\end{ex}

% Câu 22
\begin{ex}
  Một bể cá hình hộp chữ nhật có chiều dài $24\text{ dm}$, chiều rộng $2\text{ dm}$ và mực nước cao $1\text{ dm}$. Thể tích nước trong bể là bao nhiêu lít (biết $1\text{ dm}^3 = 1\text{ l}$)?
  \shortans{48}
  \loigiai{
    Thể tích là: 24 x 2 x 1 = 48 dm3 = 48 lít.
  }
\end{ex}

% Câu 23
\begin{ex}
  Một bể cá hình hộp chữ nhật có chiều dài $25\text{ dm}$, chiều rộng $2\text{ dm}$ và mực nước cao $1\text{ dm}$. Thể tích nước trong bể là bao nhiêu lít (biết $1\text{ dm}^3 = 1\text{ l}$)?
  \shortans{50}
  \loigiai{
    Thể tích là: 25 x 2 x 1 = 50 dm3 = 50 lít.
  }
\end{ex}

% Câu 24
\begin{ex}
  Một bể cá hình hộp chữ nhật có chiều dài $26\text{ dm}$, chiều rộng $2\text{ dm}$ và mực nước cao $1\text{ dm}$. Thể tích nước trong bể là bao nhiêu lít (biết $1\text{ dm}^3 = 1\text{ l}$)?
  \shortans{52}
  \loigiai{
    Thể tích là: 26 x 2 x 1 = 52 dm3 = 52 lít.
  }
\end{ex}

% Câu 25
\begin{ex}
  Một bể cá hình hộp chữ nhật có chiều dài $27\text{ dm}$, chiều rộng $2\text{ dm}$ và mực nước cao $1\text{ dm}$. Thể tích nước trong bể là bao nhiêu lít (biết $1\text{ dm}^3 = 1\text{ l}$)?
  \shortans{54}
  \loigiai{
    Thể tích là: 27 x 2 x 1 = 54 dm3 = 54 lít.
  }
\end{ex}
